\subsection{Public Key Cryptography}

\subsubsection{Key Exchange}

\content{
        Up to this point every method of sending an encrypted message requires a
        secret key known to both Alice and Bob. A big problem in cryptography is
        the following: 

        Suppose that Alice and Bob are unable to meet in person, they have to
        email each other to agree on the secret key.  Unfortunatly, Eve has
        access to Alice's emails, and thus can read any emails that are sent
        between Alice and Bob. How can Alice and Bob decide on a secret key
        without the key falling into Eve's hands? 
}{% presentation
\newpage
}{% nopresentation
\vspace{3in}
}{% comments
We need some `one-way function' that is easy to compute, but hard to undo. 

Illustrate how this can be done with the paint mixing example. In this example,
the secret key is the final paint color. 

Of course, since Alice and Bob are using computers, the one-way function will
need to be a numerical function. 
}

\subsubsection{Modular Arithmetic}

\content{
        What is modular arithmetic? 
}{% presentation
\vspace{3in}
}{% nopresentation
\vspace{3in}
}{% comments
Do clock example. $8+5 \text{"="} 1$, we will write this as $8+5\equiv 1 \mod{12}$.
}

\content{
        \begin{definition} The \textbf{modulus} is another name for the
        remainder after division.
        \end{definition}
}{% presentation
\vspace{2in}
}{% nopresentation
\vspace{2in}
}{% comments
}

\content{
        \begin{example} Compute 
        \begin{enumerate}
        \item $ 15 \mod{7}$
        \item $ 125 \mod{43}$
        \item $ 5 \mod{20}$
        \item $ -12 \mod{7}$
        \item $ 2^7 \mod{5}$
        \end{enumerate}
        \end{example}
}{% presentation
\vspace{3in}
}{% nopresentation
\vspace{3in}
}{% comments
Note that will always write our final answer as a natural number between 0 and
$n-1$.
}

\content{
        \begin{example} Your turn! Compute the following
        \begin{enumerate}
        \item $ 123 \mod{19} $
        \item $ 65 \mod{12} $
        \item $ -923 \mod{134} $
        \item $ -124 \mod{543} $
        \item $ 1045 \mod{521} $
        \item $ 8324 \mod{358}$
        \end{enumerate}
        \end{example}
}{% presentation
\vspace{3in}
}{% nopresentation
\vspace{3in}
}{% comments
}

